\chapter{Curriculum, Access, and Translation}
\label{ch:12}

Whether an education can be universal depends on more than the capacity to teach. It depends on whether instruction reaches persons in a language they understand, in a form they can use, at a cost they can bear, and with a content that acknowledges their circumstances. The failures of access that remain in the world are largely of these kinds. Children in many regions are taught in a language they do not speak at home, adults who have never been taught to read have no way to use text-based services, persons with sensory or cognitive disabilities encounter materials not designed for them, and households without reliable electricity or connectivity cannot rely on any service that assumes both. This chapter examines what capable systems can contribute to these problems, what they cannot, and what distinguishes a program of access that expands human choice from one that merely distributes a product.

Translation is the most obvious contribution and the one for which the evidence is most developed. Machine translation of text has improved to the point that for widely used language pairs it is serviceable for many practical purposes, and for languages with large digital corpora the quality approaches what a competent bilingual reader would accept in non-critical settings. The qualification, which the interview presses in its discussion of the dependence of language technology on the texts available, is that the quality is distributed very unequally. Of the several thousand languages in the world a small number account for nearly all of the data on which systems are trained, and for the rest the systems are poor or unavailable. This asymmetry reproduces the asymmetry of power that produced it: the languages with the most data are those of the societies with the most institutions, and the speakers of the others receive the least benefit and bear the greatest risk of being represented incorrectly. The program accordingly distinguishes between translation as a convenience for speakers of well-resourced languages and translation as an instrument of inclusion for speakers of others, and it assigns to the latter a priority which the market does not.

The priority is not achieved by simply collecting data from the communities concerned. The collection of the speech and writing of a minority community for the training of a commercial system raises questions of consent, ownership, and benefit that have been the subject of sustained criticism, in particular where a language is embedded in cultural practices that the community regards as not to be shared. A defensible approach treats the community as a party to the arrangement with authority over the use of its language data, grants that authority in a form that persists, ensures that the products made from the data are available to the community on terms it finds acceptable, and supports the community's own capacity to build and maintain the tools. Such an approach is slower and more expensive than extraction, and the program accepts the cost on the ground that a service built on the dispossession of its beneficiaries is not an instrument of access in the sense the book intends.

A second question concerns curriculum. The content of education is not a neutral selection from a body of settled knowledge, and universal education requires decisions about what is taught first, what is taught to all, and what is left to specialization. The mathematics of a stylized model helps to show the structure of these decisions without settling them. If the frequency of words in a language declines roughly in inverse proportion to their rank, as the empirical regularity associated with Zipf~\cite{zipf1949} describes, then a small vocabulary covers a large share of ordinary text, and the share gained by each additional word declines steadily. In the formalization, a vocabulary of one thousand words in a language of one hundred thousand covers roughly sixty per cent of running text under this stylization, and doubling it adds only about six percentage points. The lesson, which holds more generally for skills, is that a basic competence covering the bulk of everyday demands can be taught to everyone quickly, that the remaining demands are served at ever higher cost, and that the point at which universal provision should stop and specialization begin is a judgment about benefit and cost which a community must make. A capable system can reduce the cost of moving beyond the basic level for an individual who wants to, and it can personalize the order and pace of learning; it cannot decide what the community owes every child.

Adaptive pacing deserves discussion, because it is both the great promise of automated instruction and a source of subtle difficulty. A system that adjusts the difficulty of material to the performance of the student can keep each student in a zone where tasks are demanding but achievable, and the advantage over a class paced to the average is evident. The difficulty is that the adjustment is governed by a measure of performance, and the measure is not the same thing as the educational aim. A system optimized to maximize correct answers will tend to select easy material, and a system optimized to maximize engagement will tend to select material that holds attention, which is not always the same as material that educates. The remedy is to define the objective in terms of learning measured independently of the system's own measurements, to audit the behavior of the system for drift toward easy or entertaining tasks, and to allow teachers to inspect and override the sequence. It is a specific application of the general principle that the objective of an optimizing process should be set by those accountable for it and checked against an independent standard.

Access in the physical sense must also be provided for, and here the engineering of delivery matters as much as the sophistication of the software. A system that requires continuous high-bandwidth connectivity excludes the households that most need it. A design suited to the real distribution of circumstances stores content locally, operates on inexpensive hardware, tolerates intermittent power and connection, synchronizes when it can, and degrades gracefully when it cannot. The design should also provide for shared use, since in many settings a device is used by a family or a classroom and not by an individual. Provision for persons with disabilities is likewise a matter of design from the beginning and not an adjustment at the end: captions, audio description, alternative input methods, and content adapted to differences in cognitive style are requirements which, once met, tend to benefit all users. The program's measure of success is the proportion of the target population that can use the service in practice, assessed by sampling among those least likely to be reached, and not the number of installations or the volume of traffic.

The last matter is the relation between universal education and the maintenance of a public culture. A curriculum delivered through individualized systems risks replacing the common body of knowledge and the shared experiences that allow citizens to deliberate together with a set of private educations that do not overlap. The risk is real in proportion to the degree of personalization. It can be addressed by preserving a common core, taught and discussed in groups, in which students encounter one another's views, and by reserving personalization for the pace and manner of learning and not for the content of the shared core. The reasons for this are those of the earlier chapters: a society directs itself only if its members can understand one another's claims, and an education that fragments the basis for mutual understanding weakens the capacity it was meant to build.

\section*{Formalization}

Let a language have vocabulary of size $V$ and suppose the relative frequency of the word of rank $r$ is $f_r=r^{-1}/H_V$, where $H_n=\sum_{k=1}^n k^{-1}$ is the $n$-th harmonic number. The \emph{coverage} of the $K$ most frequent words is $C(K)=H_K/H_V$.

\begin{proposition}
$C$ is increasing and concave in $K$ in the sense that the marginal coverage $C(K+1)-C(K)=\bigl[H_V(K+1)\bigr]^{-1}$ is decreasing in $K$. For large $K$ and $V$, $C(K)\approx(\ln K+\gamma)/(\ln V+\gamma)$, where $\gamma\approx0.5772$ is Euler's constant, so that to double the number of words multiplies the uncovered share by approximately $1-\ln 2/(\ln V+\gamma)$ and adds $\ln2/(\ln V+\gamma)$ to coverage independent of $K$.
\end{proposition}

\begin{proof}
The marginal coverage follows from $H_{K+1}-H_K=(K+1)^{-1}$. The asymptotic formula follows from $H_n=\ln n+\gamma+O(1/n)$, and the increment for doubling is $[\ln(2K)-\ln K]/(\ln V+\gamma)=\ln2/(\ln V+\gamma)$.
\end{proof}

For $V=10^5$ one finds $C(100)\approx0.429$, $C(500)\approx0.562$, $C(10^3)\approx0.619$, $C(2\times10^3)\approx0.676$, and $C(5\times10^3)\approx0.752$. Each doubling adds about $\ln2/(\ln10^5+0.577)\approx0.057$ to coverage, in agreement with the increment from $10^3$ to $2\times10^3$ words. The pure Zipf form is a stylization, and real corpora depart from it, particularly in the tail, but the qualitative conclusion that a modest vocabulary yields the bulk of coverage and that further coverage becomes steadily more expensive is robust. It supports the staging of curricula from a universal core to optional extensions, with the stopping point of the core fixed by a collective decision that weighs the marginal benefit $[H_V(K+1)]^{-1}$ against the marginal cost of teaching the $(K+1)$-th item.

A second computation addresses the evaluation of adaptive systems. Let a student's mastery after a sequence of tasks be $m_T$ and let a system choose task difficulties $d_t$ to maximize the proportion of correct responses $\sum_t\mathbf 1[\text{correct}_t]$. If the probability of a correct response is a decreasing function $p(d_t-m_t)$ of the gap between difficulty and mastery, while the gain in mastery from a task is $g(d_t-m_t)$ with $g$ maximized at a positive gap $\delta^\ast>0$, then the proportion-correct objective is maximized by the easiest available tasks (the smallest admissible $d_t$), while the gain objective is maximized at $d_t=m_t+\delta^\ast$. The two objectives agree only if $g$ is itself maximized at the easiest task. This elementary observation is the formal reason for the requirement that the objective be defined as independently measured learning, since otherwise the optimization drives the system to the task level at which the student learns least.
